\documentclass{article}
\usepackage[T5]{fontenc}
\usepackage[utf8]{inputenc}
\usepackage[vietnam]{babel}
\usepackage{amsmath}
\usepackage{graphicx}
\usepackage{hyperref}
\usepackage{listings}
\usepackage{color}

\definecolor{codegreen}{rgb}{0,0.6,0}
\definecolor{backcolour}{rgb}{0.95,0.95,0.92}

\lstset{
    backgroundcolor=\color{backcolour},   
    commentstyle=\color{codegreen},
    basicstyle=\ttfamily\footnotesize,
    breaklines=true,                 
    numbers=left,                    
    numbersep=5pt,                  
    showstringspaces=false,
    tabsize=2
}

\title{Giải thích LAB 1: Machine Learning Project (End-to-End)}
\author{Gemini CLI Assistant}
\date{May 2026}

\begin{document}
\maketitle

\section{Mục tiêu}
Thực hiện một dự án Học máy hoàn chỉnh từ bước thu thập dữ liệu thô đến việc lưu trữ mô hình sẵn sàng sử dụng (\texttt{model.pkl}).

\section{Triển khai chi tiết}
Chúng tôi sử dụng bộ dữ liệu \textbf{California Housing} để dự đoán giá nhà. Quy trình bao gồm:

\subsection{Tiền xử lý dữ liệu (Preprocessing)}
\begin{itemize}
    \item \textbf{Xử lý giá trị thiếu:} Sử dụng \texttt{SimpleImputer} với chiến lược \texttt{median}. Đây là bước quan trọng vì hầu hết các thuật toán ML không thể làm việc với dữ liệu rỗng.
    \item \textbf{Mã hóa (Encoding):} Chuyển đổi \texttt{ocean\_proximity} (dạng chữ) sang dạng vector số bằng \texttt{OneHotEncoder}.
    \item \textbf{Chuẩn hóa (Scaling):} Sử dụng \texttt{StandardScaler}. Các thuật toán dựa trên khoảng cách hoặc tối ưu hóa (như Gradient Descent) sẽ hoạt động kém nếu các đặc trưng có thang đo quá khác biệt.
\end{itemize}

\subsection{Lựa chọn và Huấn luyện mô hình}
\begin{itemize}
    \item \textbf{Linear Regression:} Mô hình cơ bản, dễ hiểu nhưng có thể bị Underfitting nếu dữ liệu phức tạp.
    \item \textbf{Decision Tree:} Có khả năng học các mối quan hệ phi tuyến nhưng dễ bị Overfitting.
    \item \textbf{Random Forest:} Một phương pháp Ensemble (kết hợp nhiều cây quyết định) giúp cải thiện độ chính xác và giảm Overfitting.
\end{itemize}

\section{Giải thích Code chi tiết}
\begin{itemize}
    \item \textbf{Pipeline:} Chúng tôi sử dụng \texttt{sklearn.pipeline.Pipeline}. Đây là công cụ giúp tự động hóa quy trình: dữ liệu thô $\rightarrow$ xử lý số $\rightarrow$ chuẩn hóa $\rightarrow$ mô hình. Nó giúp tránh rò rỉ dữ liệu (data leakage).
    \item \textbf{fit\_transform():} Phương pháp này kết hợp \texttt{fit} (học các tham số như trung bình, độ lệch chuẩn) và \texttt{transform} (áp dụng các tham số đó để biến đổi dữ liệu). Lưu ý: Chỉ dùng \texttt{fit} trên tập huấn luyện, không bao giờ dùng trên tập kiểm tra.
    \item \textbf{GridSearchCV:} Hàm này chạy thử mô hình với tất cả các tổ hợp tham số bạn cung cấp (ví dụ: \texttt{n\_estimators} là 10, 30, 50) và sử dụng Cross-Validation để tìm ra bộ tốt nhất.
    \item \textbf{joblib.dump():} Dùng để "đóng băng" mô hình đã huấn luyện vào một file nhị phân. Khi cần dùng lại, bạn chỉ cần \texttt{load()} mà không phải huấn luyện lại từ đầu.
\end{itemize}

\section{Tối ưu hóa và Lưu trữ}
\begin{itemize}
    \item Sử dụng \texttt{GridSearchCV} để tìm bộ tham số tốt nhất cho Random Forest.
    \item Mô hình cuối cùng được lưu dưới dạng \texttt{model.pkl} bằng thư viện \texttt{joblib}.
\end{itemize}

\end{document}
